\documentclass[11pt]{article}
\usepackage[utf8]{inputenc}
\usepackage[english]{babel}
\usepackage{geometry}
\usepackage{amsmath}
\usepackage{booktabs}
\usepackage{xcolor}
\usepackage{mdframed}

\geometry{a4paper, margin=1in}

\title{Technical Specifications: OpenSourceSDRLab 7020-AD9361 (Advanced PlutoSDR)}
\author{Radiofrequency Engineering}
\date{\today}

\newmdenv[
  backgroundcolor=red!8,
  linecolor=red!60,
  linewidth=1pt,
  roundcorner=4pt,
  innertopmargin=6pt,
  innerbottommargin=6pt
]{warningbox}

\newmdenv[
  backgroundcolor=yellow!10,
  linecolor=yellow!60!black,
  linewidth=1pt,
  roundcorner=4pt,
  innertopmargin=6pt,
  innerbottommargin=6pt
]{pendingbox}

\begin{document}

\maketitle
\tableofcontents
\newpage

% =========================================================
\section{Gains}
% =========================================================

Managed through the \texttt{libiio} API (Linux Industrial I/O), which directly controls the AD9361 radio chains and the additional hardware integrated on this board:

\begin{itemize}
    \item \textbf{RX Gain (AD9361):} Range from 0 to 73 dB. Can operate in MGC (Manual Gain Control) or through AGC (Automatic) with Fast, Slow, or Hybrid profiles. Affects the integrated LNA, mixer, and TIA filters.
    \item \textbf{TX Attenuation (AD9361):} Digitally adjustable attenuation from 0 to -89.75 dB in fine 0.25 dB steps.
    \item \textbf{Integrated PA (External Hardware):} Unlike the official PlutoSDR, this board includes a power amplifier (PA) in the transmit chain.
    \begin{pendingbox}
    Direct on/off control of the PA depends on the GPIO pin configuration enabled in the modified firmware. By default, it significantly increases output power compared with the native $\approx$0 dBm of the AD9361.
    \end{pendingbox}
\end{itemize}

% =========================================================
\section{Hardware Filters and Processing}
% =========================================================

\begin{itemize}
    \item \textbf{Analog Processing (AD9361):} Transimpedance-amplifier (TIA) and baseband low-pass filters (LPF), highly configurable in real time to prevent aliasing and limit the analog channel bandwidth (from 200 kHz to 56 MHz).
    \item \textbf{Integrated Digital Filtering (AD9361):} 128-tap FIR filter blocks and CIC filters for interpolation/decimation before sending samples to the SoC.
    \item \textbf{SoC and FPGA (Xilinx Zynq XC7Z020-CLG400):} This is the largest jump compared with the PlutoSDR (which uses the 7010). The 7020 offers $\approx$3 times more logic cells (85k vs 28k), allowing complex custom bitstreams to be loaded (e.g., complete modems in VHDL/Verilog) and allowing the board to emulate a ZEDBOARD + FMCOMMS2/3 environment.
\end{itemize}

% =========================================================
\section{Frequency Ranges}
% =========================================================

\begin{itemize}
    \item \textbf{Operating Range:} 70 MHz to 6.0 GHz.
    \item \textbf{Instantaneous Bandwidth:} Up to 56 MHz.
    \item \textbf{Hardware Note:} The board uses an AD9361 transceiver (or an AD9363 officially software-unlocked to behave like a 9361), which removes the calibration bottlenecks of low-cost commercial versions that operate only between 325 MHz and 3.8 GHz.
\end{itemize}

% =========================================================
\section{Sample Rate and Resolution}
% =========================================================

Driven by the AD9361 internal converters and the upgrade to \textbf{1 GB of DDR3 RAM} (twice that of the original PlutoSDR).

\begin{table}[h]
    \centering
    \small
    \begin{tabular}{@{}p{0.24\textwidth}p{0.22\textwidth}p{0.20\textwidth}p{0.22\textwidth}@{}}
        \toprule
        \textbf{Architecture} & \textbf{ADC/DAC Resolution} & \textbf{Max. Sample Rate} & \textbf{Host Interfaces} \\
        \midrule
        2T2R (dual MIMO) & 12-bit & 61.44 MS/s & USB 2.0 / Gigabit Ethernet \\
        \bottomrule
    \end{tabular}
\end{table}

\textbf{Data-throughput note:} The standard PlutoSDR saturates quickly over USB 2.0 (practical limit $\approx$15 MS/s). Using the \textbf{Gigabit Ethernet} interface included on this board allows much higher sample rates to be sustained in network mode, approaching the RF chip limit.

% =========================================================
\section{IQ \& DC Corrections}
% =========================================================

\begin{itemize}
    \item \textbf{Internal Calibration:} The AD9361 transceiver has internal state machines that perform automatic DC offset tracking and correction (Tx and Rx) and quadrature imbalance correction (QEC).
    \item \textbf{BIST Phase (Built-In Self Test):} Calibrations run during initialization and can be forced to rerun through the \texttt{libiio} API in case of drastic temperature changes.
\end{itemize}

% =========================================================
\section{Connectors and Inputs}
% =========================================================

\begin{table}[h]
    \centering
    \small
    \begin{tabular}{p{0.18\textwidth}p{0.26\textwidth}p{0.44\textwidth}}
        \toprule
        \textbf{Connector} & \textbf{Function} & \textbf{Specification} \\
        \midrule
        Gigabit Ethernet & Network and high-speed control & RJ45. Ideal for heavy continuous streaming. (Warning: some firmware images disable USB when Ethernet is used). \\
        Micro-USB / Type-C & Power and OTG & Classic connection for gadget mode (RNDIS) or USB peripherals. \\
        RF (SMA) $\times 4$ & TX1, TX2, RX1, RX2 & 50 $\Omega$. Enables simultaneous $2\times2$ MIMO topologies. \\
        REF CLK (SMA) & Reference Clock & Input for injecting high-precision 10 MHz or 40 MHz (GPSDO). \\
        LO IN (SMA) & External Local Oscillator & Phase input for strict coherence between multiple boards. \\
        JTAG & FPGA debugging / UART & Explicitly exposed. Includes serial lines for the boot console. \\
        \bottomrule
    \end{tabular}
\end{table}

% =========================================================
\section{Safe Operating Limits}
% =========================================================

\begin{warningbox}
\textbf{Critical warnings - front end and hardware modifications}
\end{warningbox}

\vspace{4pt}
\begin{itemize}
    \item \textbf{Maximum RX power:} Do not exceed the theoretical \textbf{+2.5 dBm} (recommended 0 dBm maximum) on the RX1/RX2 ports to protect the AD9361 LNAs.
    \item \textbf{Closed TX-RX loop (Loopback):} Since this board \textbf{includes a PA in transmission}, the TX port will send much more energy than a conventional Pluto. Use attenuators of at least 40--50 dB in direct tests.
    \item \textbf{Reference Voltages:} Ensure that signals injected into REF CLK and LO IN meet the tolerable voltage levels (typically 1.8 V logic or capacitively coupled).
    \item \textbf{Heat Dissipation:} The Zynq 7020 and AD9361 generate significant heat, especially under MIMO and high sample rates. Requires a heat sink and aluminum enclosure or forced ventilation.
\end{itemize}

% =========================================================
\section{Clock and Advanced Synchronization (MIMO)}
% =========================================================

\begin{itemize}
    \item Unlike the original basic design, this OpenSourceSDRLab board introduces \textbf{explicit REF CLK and LO (Local Oscillator) inputs}.
    \item \textbf{Phase Coherence:} Connecting the same Local Oscillator (LO) to several 7020-AD9361 boards allows them to be synchronized to form larger-scale antenna arrays ($4\times4$, $8\times8$), which is fundamental in phased-array radar or beamforming research (for which MATLAB or GNU Radio is used).
\end{itemize}

% =========================================================
\section{Indicator LEDs}
% =========================================================

\begin{table}[h]
    \centering
    \begin{tabular}{p{0.12\textwidth}p{0.36\textwidth}p{0.40\textwidth}}
        \toprule
        \textbf{LED} & \textbf{Function} & \textbf{Normal state} \\
        \midrule
        \textbf{PWR} & Power & Steady after powering on the unit. \\
        \textbf{DONE} & FPGA Boot Status & Turns on when the Zynq correctly loads the firmware bitstream. \\
        \textbf{TX/RX} & Radio activity & Blinking/steady to indicate operation of the AD9361 ports. \\
        \textbf{ETH Link} & Gigabit port & Standard RJ45 signaling (Link/Act). \\
        \bottomrule
    \end{tabular}
\end{table}

% =========================================================
\section{Expansion Interfaces (JTAG and GPIO)}
% =========================================================

\subsection*{Multipurpose JTAG Port}
The manufacturer explicitly exposed this port to facilitate ``bare-metal'' development (without Linux) or to accelerate HDL simulations:
\begin{itemize}
    \item Allows code to be loaded onto the dual-core ARM Cortex-A9 or the FPGA without going through the operating system.
    \item Includes pins for the UART console, used for U-Boot diagnostics during boot and IP configuration.
\end{itemize}

\subsection*{GPIO Expansion}
Headers have been introduced with multiple GPIO pins connected directly to the FPGA logic (XC7Z020 pin bank). Ideal for:
\begin{itemize}
    \item Controlling additional external amplifiers.
    \item Synchronizing relays and transmit/receive switches (T/R switches).
    \item Hardware triggers to/from other laboratory equipment.
\end{itemize}

\end{document}
